\documentclass{article}
\usepackage{amsmath}
\usepackage{cancel}

\title{d'Alembert's Principle}

\begin{document}
\maketitle

The typical route to derive the Euler-Lagrange equations consists of stating a
new ``principle'' of physics (Hamilton's principle) and then using the calculus
of variations to derive the equations of motion. For the first pass through, it
can seem as though quantities such as the Lagrangian $\mathcal{L}$ are pulled
out of thin air. Only after working through many examples in full detail do you
begin to accept the equivalence of this new method to the more familiar Newton's
Laws.

Here we will instead start with Newton's laws and present a different principle
that will allow us to derive the Euler-Lagrange equations without presupposing
the mysterious Lagrangian. Then, having convinced ourselves of the importance of
the quantity $\mathcal{L}$ we will explore Hamilton's principle as it has many
far reaching applications across seemingly disparate topics in physics.

To begin, we define a \textit{virtual displacement} as an infinitesimal change
of the coordinates that is consistent with whatever constraints we have imposed
\emph{at a given instant of time}. We use virtual to distinguish from an actual
displacement of the system which instead occurs during an interval $dt$.

For a system in equilibrium, we expect that the total force on each $i^{th}$
particle vanishes, i.e. $\vec{F}_i = 0$ so that the virtual work
\begin{equation}
  \sum_i \vec{F}_i \cdot \delta \vec{r}_i = 0
\end{equation}
is also zero. We can decompose the net work into an applied component and a
component due to the constraints. That is,
\begin{equation}
  \vec{F}_i \equiv \vec{F}_i^{(a)} + \vec{f}_i
\end{equation}
Because forces of constraint $\vec{f}_i$ are always perpendicular to the surface
while displacements are tangential, it is reasonable to restrict ourselves to
situations in which the \emph{forces of constraint do no work}. For forces like
sliding friction, we can not make this strong assumption and would have to amend our
theory. Note however that \emph{rolling friction does no work}.

Under this assumption, we now have
\begin{equation}
  \sum_i \vec{F}_i^{(a)}\cdot \delta \vec{r}_i = 0 
\end{equation}
Now if we write Newton's second law in the suggestive form
\begin{equation}
  \vec{F}_i - \frac{d}{dt}\vec{p}_i = 0 
\end{equation}
which we can now turn into:
\begin{equation}
  \require{cancel}
  \sum_i \left( \vec{F}_i^{(a)}+\cancel{\vec{f}_i}-\frac{d}{dt}\vec{p}_i \right) = 0 .
\end{equation}
Note that I have canceled out the constraint force according to our assumption.
This equation is commonly called \textit{d'Alembert's principle}. 
It remains to turn this equation into something a little more useful. To do
that, let's expand the terms in terms of our coordinates $\{q_i\}$. In general,
we have
\begin{equation}
  \vec{r}_i = \vec{r}_i\left( q_1, q_2, ... , q_n , t \right)
\end{equation}
therefore, we have that the total differential is
\begin{equation}
  d\vec{r}_i = \sum_k \frac{\partial \vec{r}_i}{\partial q_k}dq_k + \frac{\partial \vec{r}_i}{\partial t}dt
\end{equation}
which let's us write the velocity as
\begin{equation}
  \vec{v}_i = \sum_k \frac{\partial \vec{r}_i}{\partial q_k}\dot{q_k} + \frac{\partial \vec{r}_i}{\partial t}
\end{equation}

We can define the virtual displacements to be the total differential without the
time bit. That is,
\begin{equation}
  \delta\vec{r}_i = \sum_k \frac{\partial \vec{r}_i}{\partial q_k} \delta q_k 
\end{equation}

Now we want to expand both parts of d'Alembert's equation. Starting with
the applied forces, we have
\begin{align}
  \sum_i \vec{F}_i^{(a)}\cdot \delta \vec{r}_i &= \sum_i\sum_j\vec{F}_i^{(a)}\cdot
                                                 \frac{\partial \vec{r}_i}{\partial q_j}\delta q_j \\
                                               &= \sum_j Q_j\delta q_j
\end{align}
where $Q_j$ are the \textit{generalized forces}. 

Next, let's deal with the second term involving the momenta. We have 
\begin{equation}
  \sum_i \vec{p}_i \cdot \delta\vec{r}_i = \sum_i m_i\ddot{\vec{r}}_i\cdot \delta \vec{r}_i
\end{equation}
by definition of momentum. In terms of the coordinates, this is
\begin{equation}
  \sum_i\sum_j m_i\ddot{\vec{r}}_i \cdot \frac{\partial \vec{r}_i}{\partial q_j}\delta q_j
\end{equation}

The task of solving differential equations is hard, so let's try and simplify
our lives by ripping off one of the derivatives from this equation.
In doing so, we can take advantage of the product rule to write
\begin{equation}
  \sum_i m_i\ddot{\vec{r}}_i\cdot\frac{\partial \vec{r}_i}{\partial q_j} =
  \sum_i\left[ \frac{d}{dt}\left( m_i\dot{\vec{r}_i}\cdot\frac{\partial \vec{r}_i}{\partial q_j} \right)
  - m_i\dot{\vec{r}}_i\cdot \frac{d}{dt}\left( \frac{\partial \vec{r}_i}{\partial q_j} \right)\right].
\end{equation}
Exchanging the time derivative with the partial derivative yields
\begin{equation}
  \frac{d}{dt}\left( \frac{\partial \vec{r}_i}{\partial q_j} \right) = \frac{\partial}{\partial q_j}\frac{d\vec{r}_i}{dt} = \frac{\partial \vec{v}_i}{\partial q_j}
\end{equation}
we can also see from equation for the velocity that
\begin{equation}
  \frac{\partial \vec{v}_i}{\partial \dot{{q}}_j} = \frac{\partial \vec{r}_i}{\partial q_j}
\end{equation}
so that our momentum piece can be re-written
\begin{equation}
  \sum_i m_i\ddot{\vec{r}}_i\cdot\frac{\partial \vec{r}_i}{\partial q_j} =
  \sum_i\left[ \frac{d}{dt}\left( m_i\vec{v}_i\cdot \frac{\partial \vec{v}_i}{\partial \dot{{q}}_j} \right)
  - m_i\vec{v}_i\cdot\frac{\partial \vec{v}_i}{\partial q_j}\right]
\end{equation}

Sticking it all together yields
\begin{equation}
  \sum_i\left( \vec{F}_i^{(a)}-\frac{d}{dt}\vec{p}_i \right) =
  \sum_j Q_j\delta_j - \sum_j   \sum_i\left[ \frac{d}{dt}\left( m_i\vec{v}_i\cdot
      \frac{\partial \vec{v}_i}{\partial \dot{{q}}_j} \right)
    - m_i\vec{v}_i\cdot\frac{\partial \vec{v}_i}{\partial q_j}\right]\delta q_j = 0 
\end{equation}

Let's use the chain rule to manipulate the summands. 
\begin{align}
  \frac{\partial}{\partial \dot{q}_j}\left(m_i \vec{v}_i^2) \right) &= 2m_i\vec{v}_i\cdot\frac{\partial \vec{v}_i}{\partial \dot{q}_j} \\
  \Rightarrow m_i\vec{v}_i\cdot\frac{\partial \vec{v}_i}{\partial\dot{q}_j} &= \frac{\partial}{\partial \dot{q}_j}\left( \frac{1}{2}m_i\vec{v}_i^2 \right)
\end{align}
By the same argument, we also have
\begin{equation}
  m_i\vec{v}_i \cdot\frac{\partial\vec{v}_i}{\partial q_j} = \frac{\partial}{\partial q_j}\left( \frac{1}{2}m_i\vec{v}_i^2 \right)
\end{equation}
So, because differentiation is linear, we can move the time derivative outside
of the sum to take advantage of these two relations. This results in
\begin{equation}
  \sum_j\left\{\frac{d}{dt}\left[ \frac{\partial}{\partial \dot{q}_j}\left( \sum_i \frac{1}{2}m_i\vec{v}_i^2  \right)\right] -\frac{\partial}{\partial q_j}\left(\sum_i\frac{1}{2}m_i\vec{v}_i^2  \right)-Q_j \right\} \delta q_j = 0
\end{equation}
and we know what $\sum_i\dfrac{1}{2}m_i\vec{v}_i^2$ is! It's the kinetic energy
$T$. With this substitution, we have the much simpler form
\begin{equation}
  \sum_j\left\{ \frac{d}{dt}\left( \frac{\partial T}{\partial\dot{q}_j} \right)-\frac{\partial T}{\partial q_j}-Q_j \right\}\delta q_j = 0 
\end{equation}

We would like to be able to say that the only way for this equation to hold is
for everything inside the sum other than the virtual displacements to be
identically zero. This argument is \textbf{really important}.

If our coordinates are independent of each other, then we can create an
arbitrary virtual displacement in the direction $\delta q_\alpha$ so that all
other $\delta q_j$ are zero. In this scenario, it must be that the rest of the
summand is zero for our previous equation to hold. Therefore, because the choice
of $\delta q_\alpha$ was arbitrary, we can repeat the process ad nauseum until
we are convinced that each summand must also be zero. We conclude
\begin{equation}
  \frac{d}{dt}\left( \frac{\partial T}{\partial \dot{q}_j} \right) - \frac{\partial T}{\partial q_j} - Q_j = 0
\end{equation}

We can do even better if our forces $\vec{F}_i^{(a)}$ can be derived from a
potential. Then, 
\begin{equation}
  Q_j = \sum_i \vec{F}_i^{(a)}\cdot\frac{\partial \vec{r}_i}{\partial q_j} = -\sum_i\nabla_i V\cdot\frac{\partial \vec{r}_i}{\partial q_j}
\end{equation}
Recalling connection between the geometric definition of the gradient and the
multivariable chain rule, we have
\begin{equation}
  df = \nabla\cdot d\vec{r} = \sum_i\left(\frac{\partial f}{\partial x_i} \right)dx_i
\end{equation}
so that we can write
\begin{equation}
  Q_j = -\sum_i\nabla_iV\cdot\frac{\partial\vec{r}_i}{\partial q_j} = -\frac{\partial V}{\partial q_j}
\end{equation}
Plugging this in yields
\begin{equation}
  \frac{d}{dt}\left( \frac{\partial T}{\partial \dot{q}_j} \right) - \frac{\partial(T-V)}{\partial q_j} = 0 
\end{equation}

Lastly, if the potential $V$ does not depend on the generalized velocities
$\dot{q}_j$, it must be true that the $\dot{q}_j$ derivative of $V$ vanishes
allowing us to write
\begin{equation}
  \frac{d}{dt}\left( \frac{\partial T-V}{\partial \dot{q}_j} \right) - \frac{\partial(T-V)}{\partial q_j} = 0 
\end{equation}
so that by defining the Lagrangian $\mathcal{L}=T-V$, we have the familiar
\begin{equation}
  \frac{d}{dt}\left( \frac{\mathcal{L}}{\partial \dot{q}_j} \right) - \frac{\mathcal{L}}{\partial q_j} = 0 
\end{equation}
Euler-Lagrange equations for the $j^{th}$ generalized coordinate.


In summary, we have assumed the following:
\begin{itemize}
\item Constraint forces do no work.
\item The chosen coordinates $\{q_j\}$ are independent of each other
\item The applied forces $\vec{F}_i^{(a)}$ can be derived from a scalar
  potential $V=V(\vec{r}_1,\vec{r}_2,..,\vec{r}_n, t)$
\end{itemize}


What has this accomplished? We derived the Euler-Lagrange equations of
Lagrangian dynamics without using the calculus of variations and presupposing
the importance of the Lagrangian. By simply starting with Newton's second law,
the Lagrangian formulation fell right out. There are a couple of limitations to
this method however. We were forced to assume that the coordinates were
independent of one another in order to allow us to break the summation. This
certainly is not always true-- particularly if we have any kind of constraint
function on the system $f(\vec{r}_\ell,...\vec{r}_k)=0$. This case and the issue
of velocity dependent potentials (like the ones that appear for the Lorentz
force) will require us to use different methods to derive the Euler-Lagrange
equations under different assumptions. 
\end{document}
